% Generated draft: requires mathematical and source review.
\documentclass[11pt]{article}
\usepackage[margin=25mm]{geometry}
\usepackage{fontspec}
\setmainfont{Latin Modern Roman}
\usepackage{amsmath,amssymb,mathtools,unicode-math}
\setmathfont{Latin Modern Math}
\usepackage{enumitem,needspace,xurl,hyperref}
\setlength{\parindent}{0pt}
\setlength{\parskip}{5pt}
\setlength{\emergencystretch}{3em}
\newcommand{\N}{\mathbb{N}}
\newcommand{\Z}{\mathbb{Z}}
\newcommand{\Q}{\mathbb{Q}}
\newcommand{\R}{\mathbb{R}}
\newcommand{\C}{\mathbb{C}}
\newcommand{\F}{\mathbb{F}}
\newcommand{\A}{\mathbb{A}}
\newcommand{\PP}{\mathbb{P}}
\newcommand{\E}{\mathbb{E}}
\newcommand{\eps}{\varepsilon}
\newcommand{\dd}{\,\mathrm{d}}
\begin{document}
\section*{UnsolvedMath ID: 30000987}
\textbf{Source problem code:} catalogue entry
\subsection{Definitions and mathematical statement}
All subgroups below are understood as abstract subgroups, not merely closed or open subgroups of a profinite completion. A maximal subgroup $M<G$ is a proper subgroup with no subgroup strictly between $M$ and $G$.

In the standard rooted-tree convention, a branch group is an infinite group $G$ with a faithful level-transitive action on a locally finite rooted tree such that the subgroup generated by the rigid stabilizers at each level has finite index in $G$. The rigid stabilizer at a vertex $v$ acts trivially outside the descendant subtree rooted at $v$. This definition and the intended tree convention should be compared with the report.

An infinite group is just infinite if every nontrivial normal subgroup has finite index. It is hereditarily just infinite if each of its finite-index subgroups is just infinite; this does not mean that every arbitrary subgroup is just infinite.

The two source questions are:
\[
\forall G\text{ finitely generated branch}\quad
\forall M<G\text{ maximal},\qquad [G:M]<\infty;
\]
and
\[
\forall G\text{ finitely generated hereditarily just infinite}\quad
\exists M<G\text{ maximal},\qquad [G:M]=\infty.
\]
They are retained as separate questions; the source's word ``conversely'' does not turn them into literal logical converses or identify the two classes of groups.

The supplied literature triage checked 2026-08-21 labels the bundle partially solved. It attributes to Bondarenko finitely generated branch groups with infinite-index maximal subgroups, giving a negative answer to the first universal question. It labels the hereditary-just-infinite question and a general characterization unresolved. Thus the first displayed assertion is a historical question with a reported counterexample, not a current affirmative conjecture. No precise additional characterization problem or counterexample construction is supplied in the record.
\subsection{Short English statement}
The first question asks whether all maximal subgroups of finitely generated branch groups have finite index; the source reports a negative answer. The second asks whether every finitely generated hereditarily just-infinite group has some infinite-index maximal subgroup and is reported unresolved.
\subsection{Sources}
\begin{itemize}
\item Local UnsolvedMath registry, permanent ID 30000987, source code catalogue entry; supplied complete record, labeled exact after source checking and carrying dated partial-solution triage.
\item Profinite and Asymptotic Group Theory (2009), Oberwolfach Reports, DOI 10.4171/owr/2008/28, pp. 1537--1588: \url{https://doi.org/10.4171/owr/2008/28}. Proposed-year metadata gives 2008.
\item Bondarenko counterexample: \url{https://arxiv.org/abs/1002.0320}.
\item Later branch-group study: \url{https://www.sciencedirect.com/science/article/am/pii/S0021869320303069}.
\end{itemize}
\end{document}
